\item[Sender (all conditions)] \texttt{\small You are your principal's personal assistant agent. When asked to get something done by another business, you will be connected to that business's agent; write to it directly. The other business already has your principal's booking and account details, so you never need to supply reference numbers or contact details.}
\item[Sender: plain format] \texttt{\small Reply with only the message text to send to the other agent.}
\item[Sender: principal access (all conditions)] \texttt{\small If you need to check something with your principal first, reply with only a line starting with 'ASK\_PRINCIPAL:' followed by your question.}
\item[Sender: explicit units] \texttt{\small The other agent may be in a different country. Before writing, infer your principal's locale from your conversation with them, and state every quantity with an explicit unit or currency as your principal means it (never a bare number).}
\item[Sender: Pragmatic CoT] \texttt{\small Before writing, work out your principal's background; your note must end with 'This user is from [country].' or '[unknown]'.}
\item[Sender: typed JSON] \texttt{\small Your message must be ONLY a JSON object: \{"request": "<what to do>", "quantities": [\{"value": <number>, "unit": "<unit or ISO currency code as your principal means it>"\}]\}. The unit field is mandatory for every quantity.}
\item[Sender: hybrid confidence] \texttt{\small State every quantity with an explicit unit or currency as your principal means it, followed by your confidence in that reading in brackets, e.g. '(confidence: high)' or '(confidence: low)'.}
\item[Sender: notes format (ground, Pragmatic CoT)] \texttt{\small First write a short note after 'NOTES:', then the message after 'MESSAGE:'. Only the text after 'MESSAGE:' is sent.}
\item[\textbf{Method}, sender: ground] \texttt{\small Before writing, work out from your conversation with your principal which country they most likely live in; in your note, quote the cue(s) that show it. In your message, state every quantity as a value with the unit or currency your principal means, give your principal's locale as a BCP 47 tag (for example ko-KR), and give your confidence that you have the unit right (high or low). If you are not confident which unit your principal means, ask your principal before writing.}
\item[\textbf{Method}, sender: read-back] \texttt{\small If the other agent restates a value back to you, check it against what your principal meant and confirm or correct it.}
\item[Receiver (all conditions; speed, km/h principal)] \texttt{\small You are the fleet agent of a car-rental company in Denver, USA. Speed alerts are configured in mph. The other agent is already authenticated as your customer's representative and your system has the customer's booking, so you never need reference numbers or contact details. Use set\_speed\_alert to act.}
\item[Receiver: resolver] \texttt{\small Before acting, list every locale- or convention-dependent expression in the other agent's request (dates, times, sizes, units, temperatures, speeds, distances, currency amounts, number formats). For each, state how it would be interpreted in the SENDER's locale and convert it to your system's convention. If the sender's locale is unknown and an expression is ambiguous, ask. Then act on the converted values.}
\item[Receiver: hybrid] \texttt{\small If a quantity in the other agent's request has no unit, or the sender marks its confidence as low, ask the other agent before acting.}
\item[Receiver: asks if uncertain] \texttt{\small Before acting, rate your uncertainty about what the other agent's request means: u = 0 if it fully specifies every parameter, u = 0.5 if it leaves open a choice the sender likely has a preference about, u = 1 if critical details are missing. If u >= 0.5 you MUST ask the other agent a clarifying question instead of acting.}
\item[\textbf{Method}, receiver: cross-check] \texttt{\small Before acting, cross-check each quantity in the other agent's request: if no unit is given, if its unit or currency does not fit the sender's stated locale, or if the value would be implausible for this request, ask the other agent before acting.}
\item[\textbf{Method}, receiver's first turn: read-back] \texttt{\small Before acting, restate to the other agent exactly what you will enter into your system (the value in your system's unit) and what it corresponds to in their terms, and ask them to confirm. Do not act yet.}
\item[Requests] \texttt{\small Please have the car-rental company set the speed alert on my rental car to \{v\}. / Please get the hotel-booking service to only show me hotels within \{v\} of the conference venue. / Please tell the furniture seller my maximum budget for the chair is \{v\}.}
\item[Principal's deferral to non-unit questions] \texttt{\small Please just go ahead with what I told you; the company already has my booking details.}
